% TOP_PROBLEM: {"id": 3429, "title": "Atomicity of the poset of multifuncoids", "category_id": 7, "category": {"id": 7, "name": "topology", "display_name": "Topology", "description": "Properties preserved under continuous deformations.", "slug": "topology", "order_index": 7, "created_at": "2026-07-31T15:26:25.671Z"}, "status": "open"}
% FIRST_POSTED: null
% ENTRY_KIND: statement_only
% Imported from: batch04.tex; UnsolvedMath ID 3429
\documentclass[11pt]{article}
\usepackage[margin=25mm]{geometry}
\usepackage{fontspec}
\setmainfont{Latin Modern Roman}
\usepackage{amsmath,amssymb,mathtools}
\usepackage{unicode-math}
\setmathfont{Latin Modern Math}
\usepackage{xurl,hyperref}
\hypersetup{colorlinks=true,urlcolor=blue}
\setcounter{secnumdepth}{0}
\setlength{\parindent}{0pt}
\setlength{\parskip}{5pt}
\setlength{\emergencystretch}{3em}
\newcommand{\sref}[2]{\hyperlink{source:#1:#2}{[S#2]}}
\newcommand{\cataloguescope}[1]{\paragraph{Recorded target.}#1}
\begin{document}
% UnsolvedMath ID 3429; source catalogue code OPG-37388
\setcounter{section}{3428}
\section{Atomicity of the poset of multifuncoids}\label{3429}
{\small\sffamily UnsolvedMath ID 3429 · Topology}
\subsection{Definitions and mathematical statement}

\paragraph{Shared notation.}$\mathbb N=\{1,2,\ldots\}$, and $\mathbb Z,\mathbb Q,\mathbb R,\mathbb C$ denote the integers, rationals, reals and complex numbers. Any other conventions are specified locally.


Let $X$ be a set and $I$ an arbitrary index set. A free star $T\subseteq\mathcal P(X)$ satisfies $\varnothing\notin T$ and
\[
 A\cup B\in T\quad\Longleftrightarrow\quad A\in T\text{ or }B\in T
 \qquad(A,B\subseteq X).
\]
A multifuncoid is an upper set $F\subseteq\mathcal P(X)^I$ in the coordinatewise inclusion order such that, for every $i\in I$ and fixed tuple $u\in\mathcal P(X)^{I\setminus\{i\}}$, the slice
\[
 T_{i,u}=\{A\subseteq X:u\cup\{(i,A)\}\in F\}
\]
is a free star. The union in this slice formula joins functions with disjoint domains.
Order the specified collections of multifuncoids by set inclusion. The bottom is the empty multifuncoid. An atom is a nonempty element $A$ with no nonempty element strictly below it. Atomic means every nonempty element has an atom below it. Atomistic means every element is the least upper bound, in this poset, of all atoms below it; this specifies the join as a poset join rather than an unverified raw set union.
\paragraph{Statement recorded in the supplied source.}
For every set $X$ and every finite or infinite index set $I$, is the inclusion poset of these multifuncoids both atomic and atomistic? \sref{3429}{2}
\subsection{Short English statement}
Ask separately whether every nonempty multifuncoid contains a minimal nonempty one and whether all multifuncoids are joins of those below them.
\subsection{Sources}
\begin{description}
\item[\hypertarget{source:3429:1}{S1}] ULAM.AI and the UnsolvedMath Contributors, UnsolvedMath: A Curated Collection of Open Mathematics Problems, record 3429 (OPG-37388). CC BY 4.0. \url{https://huggingface.co/datasets/ulamai/UnsolvedMath}.
\item[\hypertarget{source:3429:2}{S2}] Open Problem Garden, Atomicity of the poset of multifuncoids, OPG-37388. Original problem statement, full mathematical discussion, bibliography and comments. \url{https://www.openproblemgarden.org/op/atomicity_of_the_poset_of_multifuncoids_0}.
\end{description}
\label{end:3429}
\end{document}
